\chapter{Antenna Construction: Selected Case Studies}\label{ch:workedexample}

\section{Master Integrals Used in \textit{AntCalc}}\label{sec:masterresults}

\subsection{Three- and Four-Parton Phase-Space Masters}

\subsection{Loop and Cut-Integral Master Bases}

\subsection{Master-Integral Use Across the Implemented Routes}


\section{NLO Antenna Results}

The simplest non-trivial SMQCD antenna is $A_3^0$. As such, it may be used as an example of how
\textcolor{red}{AntCalc} handles the physics involved in building and then integrating it over the
phase space. In this chapter, we are going to go over how this antenna is mathematically computed by hand, 
and how \textcolor{red}{AntCalc} handles those same calculations on its end. 

\subsection{The Tree-Level Three-Parton Antenna $A_3^0$}\label{sec:matDerivation}

\subsubsection{Analytic Construction of the Unintegrated Antenna}

The $A_3^0$ is the tree-level final-final quark-antiquark-gluon antenna. It represents the normalised 
$\gamma^*(p)\to q(k_1)g(k_3)\bar q(k_2)$ process. This process, in the current model,
is able to produce two distinct topologies which differ by what quark line the gluon is emitted from. Those
diagrams are shown below in Figure~\ref{fig:feynDiagsA30}.

\begin{figure}[h]
    \centering
    \begin{subfigure}[b]{0.45\textwidth}
        \centering
        \begin{tikzpicture}
            \begin{feynman}
                \vertex (a) {\(\gamma^*\)};
                \vertex [right=2cm of a] (b);
                \vertex [above right=2cm of b] (q) {\(q\)};
                \vertex [above right=1cm of b] (mid);
                \vertex [below right=2cm of b] (barq) {\(\bar q\)};
                \vertex [right=1.41cm of b] (g) {\(g\)};

                \diagram* {
                (a) -- [boson] (b) -- [fermion] (q),
                (barq) -- [fermion] (b),
                (mid) -- [gluon, bend right=45] (g),
                };
            \end{feynman}
        \end{tikzpicture}
        \caption{Gluon emitted from the quark line.}
    \end{subfigure}
    \hfill
    \begin{subfigure}[b]{0.45\textwidth}
        \centering
        \begin{tikzpicture}
            \begin{feynman}
                \vertex (a) {\(\gamma^*\)};
                \vertex [right=2cm of a] (b);
                \vertex [above right=2cm of b] (q) {\(q\)};
                \vertex [below right=1cm of b] (mid);
                \vertex [below right=2cm of b] (barq) {\(\bar q\)};
                \vertex [right=1.41cm of b] (g) {\(g\)};

                \diagram* {
                (a) -- [boson] (b) -- [fermion] (q),
                (barq) -- [fermion] (b),
                (mid) -- [gluon, bend left=45] (g),
                };
            \end{feynman}
        \end{tikzpicture}
        \caption{Gluon emitted from the antiquark line.}
    \end{subfigure}
    \caption{\centering Feynman diagrams depicting the two topologies described under the $A_3^0$ antenna, corresponding to the process \(\gamma^*\to q(k_1)g(k_3)\bar q(k_2)\).}
    \label{fig:feynDiagsA30}
\end{figure}


The amplitude is given as follows directly from the diagrams,

\begin{equation}\label{eq:M30amplitude}
    \begin{split}
        \mathcal M_3^0 = e_qg_s(T^{a_3})_{i_1i_2}\left[\bar u(k_1) /\!\!\!\varepsilon^*(k_3)\frac{/\!\!\!k_1+/\!\!\!k_3}{s_{13}}/\!\!\!\varepsilon(p)\varv(k_2)-\bar u(k_1) /\!\!\!\varepsilon^*(p)\frac{/\!\!\!k_2+/\!\!\!k_3}{s_{23}}/\!\!\!\varepsilon(k_3)\varv(k_2)\right]
    \end{split}
\end{equation}
where $(T^{a_3})_{i_1i_2}$ represents the colour matrix for the quark and antiquark in positions 1 and 2 and 
gluon in position 3.


The amplitude is now squared as follows,

\begin{equation}
    \begin{split}
        |\mathcal M_3^0|^2 &= e_q^2g_s^2(T^{a_3})_{i_1i_2}(T^{a_3})_{i_1i_2}\left(M_AM_A^\dagger+M_BM_B^\dagger-M_AM_B^\dagger-M_BM_A^\dagger\right)\\
                           &\text{with } M_A = \bar u(k_1) /\!\!\!\varepsilon^*(k_3)\frac{/\!\!\!k_1+/\!\!\!k_3}{s_{13}}/\!\!\!\varepsilon(p)\varv(k_2)= \bar u(k_1) \Gamma_A\varv(k_2),\\
                           &\,\,\,\,\,\,\,\,\,\,\,\, M_B = \bar u(k_1) /\!\!\!\varepsilon^*(p)\frac{/\!\!\!k_2+/\!\!\!k_3}{s_{23}}/\!\!\!\varepsilon(k_3)\varv(k_2) = \bar u(k_1) \Gamma_B\varv(k_2).
    \end{split}
\end{equation}
which in total becomes

\begin{equation}
    \begin{split}
        = \sum_{a_3,i_1,i_2} &(T^{a_3})_{i_1i_2}(T^{a_3})_{i_1i_2}\\
                             &\sum_{s_1, s_2, \lambda_g, \lambda_\gamma}\bar u(k_1,s_1) u(k_1,s_1)\Bigg[\Gamma_A \bar\Gamma_A + \Gamma_B\bar\Gamma_B - \Gamma_A\bar\Gamma_B - \Gamma_B\bar\Gamma_A\Bigg]v(k_2,s_2)\bar v(k_2,s_2)
    \end{split}
\end{equation}

Each of the $\Gamma$ terms becomes, using Dirac Algebra \textcolor{red}{[ref to dirac algebra maths]},

\begin{equation}
    \begin{split}
        \sum_{s_1,s_2,\lambda_g,\lambda_\gamma} \bar u(k_1,s_1) u(k_1,s_1)&\left(\Gamma_A \bar\Gamma_A\right)v(k_2,s_2)\bar v(k_2,s_2) = \\
                                                                        &=\frac1{s_{13}^2}\sum_{\lambda_g,\lambda_\gamma}\varepsilon_\mu^*\varepsilon_\nu\varepsilon_\rho\varepsilon_\sigma^*\operatorname{Tr}\Big(/\!\!\!k_1\gamma^\mu(/\!\!\!k_1+/\!\!\!k_3)\gamma^\rho/\!\!\!k_2\gamma^\sigma(/\!\!\!k_1+/\!\!\!k_3)\gamma^\nu\Big)\\
                                                                        &=\frac1{s_{13}^2}g_{\mu\nu}g_{\rho\sigma}\operatorname{Tr}\Big(/\!\!\!k_1\gamma^\mu(/\!\!\!k_1+/\!\!\!k_3)\gamma^\rho/\!\!\!k_2\gamma^\sigma(/\!\!\!k_1+/\!\!\!k_3)\gamma^\nu\Big)\\
                                                                        &=\frac1{s_{13}^2}\operatorname{Tr}\Big(/\!\!\!k_1\gamma^\mu(/\!\!\!k_1+/\!\!\!k_3)\gamma^\rho/\!\!\!k_2\gamma_\rho(/\!\!\!k_1+/\!\!\!k_3)\gamma_\mu\Big)\\
                                                                        &=\left(\frac{2-d}{s_{13}}\right)^2\operatorname{Tr}(/\!\!\!k_1/\!\!\!P/\!\!\!k_2/\!\!\!P),\,\,\,/\!\!\!P = (/\!\!\!k_1+/\!\!\!k_3)\\
                                                                        &=4\left(\frac{2-d}{s_{13}}\right)^2\Big(2(k_1\cdot P)(k_2\cdot P)-(k_1\cdot k_2)P^2\Big)\\
                                                                        &=2(d-2)^2\frac{s_{23}}{s_{13}}
    \end{split}
\end{equation}

\begin{equation}
    \begin{split}
        \sum_{s_1,s_2,\lambda_g,\lambda_\gamma} \bar u(k_1,s_1) u(k_1,s_1)&\left(\Gamma_B \bar\Gamma_B\right)v(k_2,s_2)\bar v(k_2,s_2) = \\
                                                                        &=2(d-2)^2\frac{s_{13}}{s_{23}}
    \end{split}
\end{equation}

\begin{equation}
    \begin{split}
        \sum_{s_1,s_2,\lambda_g,\lambda_\gamma} \bar u(k_1,s_1) u(k_1,s_1)&\left(\Gamma_A \bar\Gamma_B\right)v(k_2,s_2)\bar v(k_2,s_2) = \\
                                                                        &=2(d-2)\frac{(d-4)s_{13}s_{23}-2s_{12}^2-2s_{12}(s_{13}+s_{23})}{s_{13}s_{23}}
    \end{split}
\end{equation}

\begin{equation}
    \begin{split}
        \sum_{s_1,s_2,\lambda_g,\lambda_\gamma} \bar u(k_1,s_1) u(k_1,s_1)&\left(\Gamma_B \bar\Gamma_A\right)v(k_2,s_2)\bar v(k_2,s_2) = \\
                                                                        &=2(d-2)\frac{(d-4)s_{13}s_{23}-2s_{12}^2-2s_{12}(s_{13}+s_{23})}{s_{13}s_{23}}
    \end{split}
\end{equation}
and computing the colour matrix sum as

\begin{equation}
    \sum_{a_3,i_1,i_2} (T^{a_3})_{i_1i_2}(T^{a_3})_{i_1i_2} = \frac12(N^2-1)
\end{equation}
we are finally able to reach the final result, as follows, using dimensional regularisation as 
$d\to 4-2\epsilon$ \textcolor{red}{[ref to dimreg]},

\begin{equation}\label{eq:M30interference}
    \begin{split}
        |\mathcal M_3^0|^2 &= e_q^2g_s^2(N^2-1)\Bigg(2(d-2)^2\left(\frac{s_{13}}{s_{23}}+\frac{s_{23}}{s_{13}}\right)-4(d-2)\frac{(d-4)s_{13}s_{23}-2s_{12}^2-2s_{12}(s_{13}+s_{23})}{s_{13}s_{23}}\Bigg)\\
                           &= 2(d-2)e_q^2g_s^2(N^2-1)\Bigg((d-2)\left(\frac{s_{13}}{s_{23}}+\frac{s_{23}}{s_{13}}\right)-2\frac{(d-4)s_{13}s_{23}-2s_{12}^2-2s_{12}(s_{13}+s_{23})}{s_{13}s_{23}}\Bigg)\\
                           &= \frac8{s_{13}s_{23}}\Bigg(2q^2s_{12}+s_{13}^2+s_{23}^2-2\epsilon\left(q^2s_{12}+s_{13}^2+s_{23}^2+s_{13}s_{23}\right)+\epsilon^2\left(s_{13}+s_{23}\right)^2\Bigg)
    \end{split}
\end{equation}

The $A_3^0$ antenna also requires the $\mathcal M_2^0$ amplitude. This amplitude corresponds to the 
$\gamma^*\to q\bar q$ Born process. The amplitude is

\begin{equation}
    \mathcal M_2^0 = e_q\delta_{i_1i_2}\bar u(k_1)/\!\!\!\varepsilon(p)\varv(k_2).
\end{equation}

Applying the same spin, polarization and colour sums as above, and considering that, for a two 
final-state particles system $s_{12}=q^2$, we get

\begin{equation}\label{eq:M20interference}
    \begin{split}
        |\mathcal M_2^0|^2 = 4e_q^2(d-2)Ns_{12} &= 4e_q^2(d-2)Nq^2\\
                                                       &= 8e_q^2(1-\epsilon)Nq^2
    \end{split}
\end{equation}


The antenna is now finally ready to be built. That is accomplished using the following formula,
\begin{equation}
    A_3^0 = \frac1{2g_s^2C_F}\frac{|\mathcal M_3^0|^2}{|\mathcal M_2^0|^2}.
\end{equation}

By taking $C_F = (N^2-1)/(2N)$, the formula yields, using Eqs.~\ref{eq:M30interference} and \ref{eq:M20interference},

\begin{equation}\label{eq:A30}
    \begin{split}
        A_3^0 &= \frac{2q^2s_{12}+s_{13}^2+s_{23}^2}{q^2s_{13}s_{23}}-\epsilon\frac{(s_{13}+s_{23})^2}{q^2s_{13}s_{23}}\\
              &= \frac1{q^2}\Bigg[\left(\frac{s_{13}}{s_{23}}+\frac{s_{23}}{s_{13}}+2\frac{s_{12}q^2}{s_{13}s_{23}}\right)-\epsilon\left(2+\frac{s_{13}}{s_{23}}+\frac{s_{23}}{s_{13}}\right)\Bigg]
    \end{split}
\end{equation}

\subsubsection{Phase-Space Integration through $R_3$}\label{sec:antIntMaths}

We now integrate the $A_3^0$ antenna over the appropriate (3 final-state particle, in this case) 
phase-space. This integral is computed in $d$ dimensions and over the phase-space measure in terms
of the Mandelstam invariants ($s_{13}$, $q^2$, etc.). The integral takes the form,

\begin{equation}
    \begin{split}
        \mathcal A_3^0 &= C(\epsilon,1)\int d\Phi_{ijk}\,A_3^0\\
                       &= \mathcal N(\epsilon,1)\,(4\pi)^{3-2d}\left(q^2\right)^{1-\frac d2}\int ds_{12}ds_{13}ds_{23}d\Omega\,(s_{12}s_{13}s_{23})^{\frac{d-4}{2}}\delta\left(q^2-s_{12}-s_{13}-s_{23}\right)\,A_3^0.
    \end{split}
\end{equation}
with $d\Phi_{ijk} = d\Phi_3/\Phi_2$ and with the derivation of the $d\Phi_3$ measure being 
present in \textcolor{red}{[ref to this derivation]}.
$C(\epsilon,1)$ is the integral normalisation which is here given as 
\begin{equation}
    C(\epsilon,1) = 8\pi^2(4\pi)^{-\epsilon}e^{\epsilon\gamma_E}.
\end{equation}

This integral is quite complex and spans across all $A_3^0$ terms. Accordingly, we are not
going to compute this integral directly: instead we are going to use IBP reduction to reduce the 
expression to a linear combination of master integrals, compute those master integrals and then 
use the result to compute the integrated antenna. This process is explained in
\textcolor{red}{[ref to IBP section]}.

After applying IBP identities and reducing the integral to a linear combination of this master
integral\footnote{This process is shown completely in \textcolor{red}{[ref to ibp reduction appendix]}.},
we reach the following expression,
\begin{equation}\label{eq:A30intermsofR3}
    \mathcal A_3^0 = \frac{C(\epsilon,1)}{\Phi_2}\frac{2(2-6\epsilon+5\epsilon^2-2\epsilon^3)}{q^2\epsilon^2}R_3
\end{equation}
This antenna has a single master: $R_3$ \textcolor{red}{[ref to master integrals section/appendix]}.
This master integral is the integral over phase-space itself,

\begin{equation}\label{eq:R3}
    R_3 = \int d\Phi_3 = \frac{2^{-7+4\epsilon}\pi^{-3+2\epsilon}\Gamma^3(1-\epsilon)}{\Gamma(2-2\epsilon)\Gamma(3-3\epsilon)}\left(q^2\right)^{1-2\epsilon}.
\end{equation}
After using the expression in Eq.~\ref{eq:R3} explicitly in Eq.~\ref{eq:A30intermsofR3}becomes
\begin{equation}\label{eq:A30withR30Sub}
    \begin{split}
        \mathcal A_3^0 &= \frac{C(\epsilon,1)}{\Phi_2}\left(q^2\right)^{-2\epsilon}(4\pi)^{-3+2\epsilon}(2-6\epsilon+5\epsilon^2-2\epsilon^3)\frac{\Gamma^3(1-\epsilon)}{\epsilon^2\Gamma(3-3\epsilon)\Gamma(2-2\epsilon)}\\
                       &= \left(q^2\right)^{-\epsilon}e^{\epsilon\gamma_E}(1-2\epsilon)\big(2+\epsilon(\epsilon+2)\big)\frac{\Gamma^2(-\epsilon)}{\Gamma(3-3\epsilon)}
    \end{split}
\end{equation}

and, after expanding the integrated expression as a Laurent series in $\epsilon$, we finally 
arrive at the infrared pole structure,

\begin{equation*}
    \begin{split}
        \mathcal A_3^0 = \left(q^2\right)^{-\epsilon}\Bigg[\frac1{\epsilon^2}+\frac3{2\epsilon}&+\frac{57-7\pi^2}{12}+\epsilon\left(\frac{84-21\pi^2-8\psi^{(2)}(1)+108\psi^{(2)}(3)}{24}\right)\\
        &+\epsilon^2\left(\frac{35640-3990\pi^2-71\pi^4-720\psi^{(2)}(1)+9720\psi^{(2)}(3)}{1440}\right)\Bigg]
    \end{split}
\end{equation*}
\begin{equation}
    = \left(q^2\right)^{-\epsilon}\Bigg[\frac1{\epsilon^2}+\frac3{2\epsilon}+\frac{19}4-\frac{7\pi^2}{12}+\epsilon\left(\frac{109}{8}-\frac{7\pi^2}8-\frac{25\zeta_3}{3}\right)+\epsilon^2\left(\frac{639}{16}-\frac{133\pi^2}{48}-\frac{71\pi^4}{1440}-\frac{25\zeta_3}{2}\right)\Bigg]
\end{equation}
where the usual polygamma identities, $\psi^{(2)}(1) = -2\zeta_3$ and $\psi^{(2)}(3) = -2\zeta_3 + \frac94$,
were employed.

\subsubsection{Implementation in \textit{AntCalc}}

Using \textcolor{red}{AntCalc}, building the antenna is very simple, since the
function \texttt{BuildAntenna[type, numFinalParticles, numLoops]} already encodes the entire 
logical chain and is able to return all intermediate steps using the appropriate options. 

First, by running the function naively we get the final result in Eq.~\ref{eq:A30},

\begin{mdframed}[backgroundcolor=gray!10, linewidth=0pt]
\texttt{In[2]:= BuildAntenna[A,3,0]//Collect[\#,Epsilon]\&}\\
\hfill
\texttt{Out[2]:=}
    \begin{equation*}
        \frac{2s_{12}}{q^2 s_{13}} + \frac{2s_{12}}{q^2 s_{23}} + \frac{2s_{12}^2}{q^2 s_{13}s_{23}} + \frac{s_{13}}{q^2 s_{23}} + \frac{s_{23}}{q^2 s_{13}} + \epsilon\left(\frac{-2}{q^2} - \frac{s_{13}}{q^2 s_{23}} - \frac{s_{23}}{q^2 s_{13}}\right)
    \end{equation*}
\end{mdframed}

We can also use the \texttt{ReturnRecord} option to access all intermediate steps, and check that 
\textcolor{red}{AntCalc} goes through the same steps as the by-hand derivation in
Sec.~\ref{sec:matDerivation}. We can see that with,

\begin{mdframed}[backgroundcolor=gray!10, linewidth=0pt]
    \texttt{In[3]:= \\BuildAntenna[A,3,0,ReturnRecord->True]["IntermediateSteps"]}\\
    \texttt{Out[3]:=}

    \texttt{Amplitude:}
    \begin{equation*}
        T^{a_3}_{i_1 i_2}\left[
        \frac{\bar{u}(k_1)\,/\!\!\!\varepsilon^*(k_3)\,(/\!\!\!k_1+/\!\!\!k_3)\,/\!\!\!\varepsilon(p)\,v(k_2)}{(k_1+k_3)^2}
        -
        \frac{\bar{u}(k_1)\,/\!\!\!\varepsilon(p)\,(/\!\!\!k_2+/\!\!\!k_3)\,/\!\!\!\varepsilon^*(k_3)\,v(k_2)}{(k_2+k_3)^2}
        \right]
    \end{equation*}

    \texttt{Interference:}
    \begin{equation*}
        \frac{4(\epsilon-1)(N^2-1)\Big(-2s_{12}^2 - 2s_{12}(s_{13}+s_{23}) + (\epsilon-1)s_{13}^2 + 2\epsilon\, s_{13}s_{23} + (\epsilon-1)s_{23}^2\Big)}{s_{13}s_{23}}
    \end{equation*}

    \texttt{Result:}
    \begin{equation*}
        -\frac{2\epsilon}{q^2} + \frac{2s_{12}^2}{q^2 s_{13}s_{23}} + \frac{2s_{12}}{q^2 s_{13}} + \frac{2s_{12}}{q^2 s_{23}} + \frac{s_{13}}{q^2 s_{23}} - \frac{\epsilon\, s_{13}}{q^2 s_{23}} + \frac{s_{23}}{q^2 s_{13}} - \frac{\epsilon\, s_{23}}{q^2 s_{13}}
    \end{equation*}
\end{mdframed}
where \texttt{Amplitude}, \texttt{Interference} and \texttt{Result} respectively refer to the
expressions in Eqs.~\ref{eq:M30amplitude}, \ref{eq:M30interference} and \ref{eq:A30}.


Integrating antennae using \textcolor{red}{AntCalc} is also simple. Every step in \ref{sec:antIntMaths}
is handled by a single function: \texttt{IntegrateAntenna[antenna]}.

Before going over how this function is used, a small clarification needs to be made: the 
\texttt{IntegrateAntenna} function is not compatible with the default output of the 
\texttt{BuildAntenna} function described earlier. This is so that the built antenna is immediately 
useful and easily readable. As described in Chapter~\ref{ch:packageframework}, however, the IBP reduction
steps are handled by the LiteRed2 package, which does not expect Mandelstam invariants, but rather,
momenta. Accordingly, to use \texttt{BuildAntenna}'s output in \texttt{IntegrateAntenna} we 
must toggle the relevant flag: \texttt{BuildAntenna[A,3,0,IntegrableForm->True]}. This
will return a large association with all required information for integrating the antenna.

We may now proceed to the integration. The whole process goes as follows,

\begin{mdframed}[backgroundcolor=gray!10, linewidth=0pt]
    \texttt{In[4]:= intA30 = BuildAntenna[A,3,0,IntegrableForm->True]};\\
    \hfill
    \texttt{In[5]:= IntegrateAntenna[intA30,ExpansionOrder->2]}\\
    \hfill
    \texttt{Out[5]:=}
        \begin{equation*}
            \frac{19}{4}+\epsilon^{-2}+\frac3{2\epsilon}-\frac{7\pi^2}{12}+\epsilon^2\left(\frac{639}{16}-\frac{133\pi^2}{48}-\frac{71\pi^4}{1440}-\frac{25\zeta_3}{2}\right)+\epsilon\left(\frac{109}{8}-\frac{7\pi^2}{8}-\frac{25\zeta_3}{3}\right)
        \end{equation*}
\end{mdframed}

Much like with \texttt{BuildAntenna}, \texttt{IntegrateAntenna} also allows us to inspect
its intermediate steps. Here, however, the available options show the master integral linear
combination expression both with $R_3$ and with it substituted by the appropriate expression, 
that is, the equivalent to Eqs.~\ref{eq:A30intermsofR3} and \ref{eq:A30withR30Sub} before
normalisation. The code and results are,

\begin{mdframed}[backgroundcolor=gray!10, linewidth=0pt]
    \texttt{In[6]:= resA30 = IntegrateAntenna[intA30,ExpansionOrder->2,\\ReturnRecord->True];}\\
    \hfill
    \texttt{In[7]:= resA30["MasterCombination"]//Simplify//Collect[\#,R3]\&}\\
    \hfill
    \texttt{Out[7]:=}
        \begin{equation*}
            \frac{-4(d-3)\big(4+d(\epsilon-2)-4\epsilon\big)}{(d-4)^2q^2}R_3
        \end{equation*}
    \texttt{In[8]:= resA30["MasterSubstitutedExpression"]//Simplify}\\
    \hfill
    \texttt{Out[8]:=}
        \begin{equation*}
            -\left(q^2\right)^{-2\epsilon}\frac{2^{-5+4\epsilon}(d-3)\big(4+d(\epsilon-2)-4\epsilon\big)\pi^{-3+2\epsilon}\Gamma^3(1-\epsilon)}{(d-4)^2\Gamma(3-3\epsilon)\Gamma(2-2\epsilon)}
        \end{equation*}
\end{mdframed}

The results here are given as a combination of $d$ and $\epsilon$. This is a result of the way
the function processes each step. Since these are snapshots of what is happening inside the
function itself, they show the expressions before dimension regularisation is applied.

Finally, \textcolor{red}{AntCalc} also enables us to forgo the work we did building and
the integrating the antenna with the use of the \texttt{BuildAndIntegrateAntenna[type,
numFinalParticles,numLoops]} function. It goes over the whole process, as we did, and outputs the 
$\epsilon$ Laurent series,

\begin{mdframed}[backgroundcolor=gray!10, linewidth=0pt]
    \texttt{In[9]:= BuildAndIntegrateAntenna[A,3,0,ExpansionOrder->2]}\\
    \hfill
    \texttt{Out[9]:=}
        \begin{equation*}
            \frac{19}{4}+\epsilon^{-2}+\frac3{2\epsilon}-\frac{7\pi^2}{12}+\epsilon^2\left(\frac{639}{16}-\frac{133\pi^2}{48}-\frac{71\pi^4}{1440}-\frac{25\zeta_3}{2}\right)+\epsilon\left(\frac{109}{8}-\frac{7\pi^2}{8}-\frac{25\zeta_3}{3}\right)
        \end{equation*}
\end{mdframed}

The agreement between the two sections in the present chapter confirms that \textcolor{red}{AntCalc}
completes the same function we would do by-hand, going through the same steps. The same
pipeline extends to the more complex antennae treated in this thesis, with their integrated 
results collected in Appendix~\ref{ap:AntennaTrees}.

\subsection{Local Validation of $A_3^0$}

\subsubsection{Soft Limit and the Eikonal Factor}

\subsubsection{Collinear Limit and the Altarelli--Parisi Splitting Function}

\subsection{The Virtual $A_2^1$ Contribution}


\section{NNLO Antenna Results}

\subsection{Two-Parton Contributions: the $A_2^2$ Set}

\subsection{Three-Parton Contributions: the $A_3^1$ Set}

\subsection{Four-Parton Contributions}

\subsubsection{The $A_4^0$ Antenna and Its Reduction}

\subsubsection{The $\tilde A_4^0$, $B_4^0$, and $C_4^0$ Results}


\section{Summary of the Massless Antenna Set}\label{sec:antennaresultsummary}


\section{Exploratory Massive Extension: $A_3^0$}\label{sec:massivea30}

\subsection{Massive Unintegrated Antenna}

\subsection{Reduction to Massive Master Integrals}

\subsection{Present Status and Limitation}